\documentclass[10pt,a4paper]{amsart}
\usepackage[margin=19mm]{geometry}
\usepackage{amsmath,amssymb,mathtools}
\usepackage[scr=rsfs,cal=euler]{mathalfa}
\usepackage{xfrac}
\usepackage[unicode,hidelinks,bookmarks=false]{hyperref}
\newcommand{\ie}{i.e.\ }
\newcommand{\cf}{cf.\ }
\newcommand{\resp}{respectively\ }
\newcommand{\Subsection}{\subsection}
\providecommand{\nobreakdash}{\nobreak\hbox{-}}
\setlength{\emergencystretch}{2em}
\multlinegap0pt
\usepackage{amssymb}
\usepackage{mathtools}
\usepackage[shortlabels]{enumitem}
\usepackage{xcolor}
\newtheorem{lemma}{Lemma}
\newcommand{\Harm}{\mathbf H}
\newcommand{\normQ}[1]{\left\lVert #1\right\rVert_Q}
\title{A Proof of Fraenkel's Conjecture}
\author{Hu Tan, Ying Zhang}
\date{Source: arXiv:2609.01570v1 (CC BY 4.0)}
\begin{document}
\maketitle

\noindent\emph{Source and licence.} A Proof of Fraenkel's Conjecture. Version arXiv:2609.01570v1 (CC BY 4.0), distributed by arXiv under the Creative Commons Attribution 4.0 International licence, http://creativecommons.org/licenses/by/4.0/, as recorded in the arXiv licence metadata for this exact version and shown on the abstract page https://arxiv.org/abs/2609.01570v1. Full licence notice: \texttt{licenses/arxiv-2609.01570-CC-BY-4.0.txt}.\par\smallskip

\noindent\textit{Editorial scope.} Counted unit: the article's lemma lem:dyadic-close (source0.tex lines 1668--1681). The article's complete original proof of it is delivered with it (source0.tex lines 1683--1847, one \texttt{proof} environment of 165 lines). No mathematics is written, repaired or re-derived: the statement and the proof are the article's own lines, in the article's own notation, and no retained line is substituted or re-flowed.  No further source-local context is needed: every cross-reference of every retained line resolves inside the two delivered spans.  Nothing was omitted from any retained span, so the package carries no omission record.  No retained line contains a citation, so the wrapper carries no bibliography and no bibliography entry.  No retained line contains a figure environment, an image-inclusion command or drawing code, so no image is delivered and the package declares no asset dependency.  Editorial formatting only, in addition: an AMS \texttt{amsart} preamble that restates the article's own declarations and macros used by the retained lines, namely the environments \texttt{lemma} on separate counters, together with the standard packages those lines need.  The retained environments print in this document as Lemma 1 in order of appearance, whereas the article prints the same items with the numbering of its own full document.  The complete native source of the article as distributed in the arXiv e-print for this exact version is bundled unchanged as \texttt{sources/209125/source0.tex}.
\begin{lemma}[Dyadic--close harmonic envelope]\label{lem:dyadic-close}
Let
\[
 H_y^{\rm good}:=\frac1y+
 \sum_{\substack{b\in A\setminus\{1,y\}\\ b\text{ is good}}}
 \frac1{\normQ{yb^{-1}}};
\]
thus the superscript includes the anchor \(b=1\) as well as every selected
good element other than \(y\).  This contribution to the \(y\)-row satisfies
\begin{equation}\label{eq:dyadic-close}
 H_y^{\rm good}
 <\frac2m+\frac{10J}{3y}+\frac{2\Harm_{y-1}}{y-1}.
\end{equation}
\end{lemma}

\begin{proof}
Because \(y\) is good, it divides one of \(Q-1,Q+1\).  Thus choose
\(\eta\in\{\pm1\}\) and an integer \(D\) with \(Q=Dy+\eta\).  Since
\(y<Q/3\), necessarily \(D\ge3\).  Put
\[
 \mathcal S_y=\{1\}\cup
 \{b\in A\setminus\{1,y\}:b\text{ is good}\},
\]
and partition \(\mathcal S_y\) into
\[
 \begin{aligned}
 P&=\{b\in\mathcal S_y:b<y,\ b\mid y\},&
 X&=\{b\in\mathcal S_y:b<y,\ b\nmid y\},\\
 \mathcal C&=\{b\in\mathcal S_y:y<b<2y\},&
 F&=\{b\in\mathcal S_y:b\ge2y\}.
 \end{aligned}
\]
By the choice of \(y\), every element of \(\mathcal S_y\) below \(y\) is
dyadic.  If \(p_0=\sum_{b\in P}b\), then
\[
 p_0\le1+2+4+\cdots+2^v=2^{v+1}-1=\frac{2y}{m}-1.
\]
For \(b\in P\), the integer \(y/b<Q/2\) satisfies
\((y/b)b=y\), so it is the least cyclic representative:
\(d_{yb}=y/b\).  Thus the \(P\)-contribution is exactly \(p_0/y\).

\smallskip
\noindent\emph{Dyadic divisors and the far shell.}
For \(b\in F\), goodness permits a sign \(\varepsilon\in\{\pm1\}\)
and a positive integer \(L\) with
\(Lb=Q+\varepsilon\).  Then
\(b^{-1}\equiv\varepsilon L\pmod Q\) and
\(d_{yb}=\normQ{Ly}\).  Since \(b\ge2y\),
\(Ly\le(Q+1)/2\).  If \(Ly\le Q/2\), then
\(bd_{yb}=y(Q+\varepsilon)\ge y(Q-1)\).  If \(Ly>Q/2\), integrality
and the oddness of \(Q\) (because \(2\in A\) is a unit) force
\(Ly=(Q+1)/2\).  Hence
\[
 b=\frac{Q+\varepsilon}{L}
   =2y\frac{Q+\varepsilon}{Q+1}.
\]
If \(\varepsilon=-1\), this is strictly smaller than \(2y\), contrary to
\(b\in F\).  Therefore \(\varepsilon=1\) and \(b=2y\); hence
\(bd_{yb}=b(Q-1)/2=y(Q-1)\) again.  Thus
\begin{equation}\label{eq:far-baseline}
 \frac1{d_{yb}}\le\frac{b}{y(Q-1)}\qquad(b\in F).
\end{equation}
The total numerator weight available to \(F\) is at most
\(Q-y-p_0\).  Using this budget and the fact that
\(p/y+(Q-y-p)/(y(Q-1))\) increases with \(p\), we obtain
\begin{equation}\label{eq:P-F}
 \sum_{b\in P\cup F}\frac1{d_{yb}}
 \le\frac{p_0}{y}+\frac{Q-y-p_0}{y(Q-1)}\le\frac2m.
\end{equation}
For the final inequality, the function
\(\phi(p)=p/y+(Q-y-p)/[y(Q-1)]\) is increasing, and at
\(p=2y/m-1\) its gap from \(2/m\) is
\[
 \frac2m-\phi\!\left(\frac{2y}{m}-1\right)
 =\frac{y+(2y/m-1)-1}{y(Q-1)}>0.
\]

\smallskip
\noindent\emph{Smaller dyadic nondivisors.}
Every \(b\in X\) is \(b=2^{v+j}\), \(1\le j\le J\).  Choose
\(\varepsilon\in\{\pm1\}\) and \(L\) with
\(Q=L2^{v+j}+\varepsilon\), and write
\(m=\rho2^j+r\), \(0<r<2^j\).  From
\[
 D2^vm+\eta=L2^{v+j}+\varepsilon
\]
we get the integer
\[
 c=\frac{\eta-\varepsilon}{2^v},\qquad
 h=\frac{Dr+c}{2^j},\qquad 1\le h\le D-1.
\]
The preceding equality shows that \(\eta-\varepsilon\) is divisible by
\(2^v\), so \(c\) is an integer.  After division by \(2^v\), it reads
\[
 L2^j=Dm+c=D\rho2^j+Dr+c,
\]
so \(h=(Dr+c)/2^j=L-D\rho\) is an integer.  Since \(|c|\le2\),
\(D\ge3\), and \(1\le r<2^j\),
\[
 Dr+c\ge D-2\ge1,
 \qquad
 Dr+c\le D(2^j-1)+2<D2^j.
\]
Thus \(1\le h\le D-1\), and \(L=D\rho+h\).  Since
\(2^{v+j}L\equiv-\varepsilon\pmod Q\), one has
\((2^{v+j})^{-1}\equiv-\varepsilon L\pmod Q\), and the distance is
\(d_{yb}=\normQ{Ly}\).  Moreover,
\[
 Ly=(D\rho+h)y=\rho(Q-\eta)+hy
 \equiv hy-\rho\eta\pmod Q.
\]
The following two integers sum to \(Q\), and hence are the complementary
representatives of this residue:
\[
 hy-\rho\eta,\qquad (D-h)y+(\rho+1)\eta.
\]
The first is at least \(y-\rho\), while the second is at least
\(y-\rho-1\).  Thus both are positive and at least \(y-\rho-1\).  Since
\(2\rho\le m-1\le y-1\), we have
\(\rho\le(y-1)/2\), and therefore
\[
 d_{yb}\ge y-\rho-1\ge\frac{y-1}{2},
 \qquad \frac1{d_{yb}}\le\frac2{y-1}<\frac{10}{3y}.
\]
There are at most \(J\) such dyadic elements, giving the middle term of
\eqref{eq:dyadic-close}.

\smallskip
\noindent\emph{The close interval.}
Finally, let \(b=y+k\in\mathcal C\), and choose
\(\varepsilon\in\{\pm1\}\), \(t\ge3\), with
\(Q=t(y+k)+\varepsilon\).  The bound \(t\ge3\) follows from
\(Q>3b=3(y+k)\): if \(\varepsilon=1\), then
\(t=(Q-1)/b>3-1/b\), and if \(\varepsilon=-1\), the bound is even
larger.  Direct reduction gives
\[
 d_{yb}=tk+\varepsilon=(D-t)y+\eta.
\]
Indeed, subtracting \(Q=t(y+k)+\varepsilon\) from
\(Q=Dy+\eta\) gives the equality of the two displayed expressions.
Moreover, \(Q-d_{yb}=ty\), and
\(ty>d_{yb}\) is equivalent to \(t(y-k)>\varepsilon\), which holds
because \(1\le k<y\) and \(t\ge3\).  Hence the displayed positive
integer is the least cyclic representative.  Put \(\ell=D-t\).  The
identity \(d_{yb}=\ell y+\eta>0\) first gives \(\ell\ge0\).  Also
\(d_{yb}=tk+\varepsilon\ge3\cdot1-1=2\), whereas \(\ell=0\) would give
\(d_{yb}=\eta\le1\).  Hence \(\ell\ge1\).
For fixed \(\ell\) and \(\varepsilon\), the integer \(t=D-\ell\) is
fixed and \(Q=t(y+k)+\varepsilon\) determines at most one \(k\), hence at
most one \(b\).  Since \(d_{yb}=\ell y+\eta\ge\ell(y-1)\), summation over
the two signs gives
\begin{equation}\label{eq:C-first}
 \sum_{b\in\mathcal C}\frac1{d_{yb}}
 \le\frac{2\Harm_{D-1}}{y-1}.
\end{equation}
For a second estimate, direct algebra gives
\[
 d_{yb}>\frac{k(D-2)}2.
\]
After multiplying the desired inequality by the positive number
\(2(y+k)\), direct substitution gives
\[
 2(y+k)\left(d_{yb}-\frac{k(D-2)}2\right)
 =Dk(y-k)+2k(y+k)+2k\eta+2\varepsilon y.
\]
Since \(\eta,\varepsilon\ge-1\), the right side is at least
\[
 Dk(y-k)+2(k-1)(y+k)>0.
\]
Therefore
\begin{equation}\label{eq:C-second}
 \sum_{b\in\mathcal C}\frac1{d_{yb}}
 <\frac{2\Harm_{y-1}}{D-2}.
\end{equation}
If \(D\le y\), then \(\Harm_{D-1}\le\Harm_{y-1}\), so use
\eqref{eq:C-first}.  If \(D\ge y+1\), then \(D-2\ge y-1\), so use
\eqref{eq:C-second}.  In either case the close contribution is at most
\(2\Harm_{y-1}/(y-1)\).  Together with
\eqref{eq:P-F}, this proves \eqref{eq:dyadic-close}.
\end{proof}

\end{document}
